\section{Boundaries}

Batch effects are expected in real election administration. Different modes,
times, geographies, ballot styles, and processing streams can all create
group-level differences. A correlation with a scanner or batch is therefore not
proof that the scanner or batch caused the pattern, and it is not proof of
fraud. It is a reason to check records, procedures, chain of custody, and
ordinary operational explanations.

The method is descriptive rather than causal. It identifies a grouping-level
association after a stated comparison, but causation belongs to a separate
investigation. A weak baseline can create false positives by comparing mail
batches to election-day precincts, early uploads to final uploads, or scanners
serving different ballot styles. The baseline population and feature set are
therefore part of the result, not background metadata.

The lead boundary governs every report: forensic analytics prioritize
investigation; they do not prove fraud, certify outcomes, or replace
risk-limiting audits. W.03 reports must remain separate from V-series evidence
and must carry their caveats with their scores.
